\tinysidebar{\debug{exercises/{test/question.tex}}}
\begin{ex}
You are given the following code that declares a 5-by-5 array of characters, initializing it with spaces.
  
\begin{consolethree}
const int ROW_SIZE = 5;

int x[ROW_SIZE][ROWS_SIZE]= {{0}};

// YOUR CODE

for (int row = 0; row < ROW_SIZE; row++)
{
     for (int col = 0; col < ROW_SIZE; col++)
     {
       std::cout << x[row][col]  << ' ';
     }
     std::cout << std::endl;
}

std::cout << std::endl; 
\end{consolethree}

Write a for-loop that sets the values on the main diagonal of the array
to \texttt{'*'}. In other words the output
should be

\begin{consolethree}
*

*

*

*

*
\end{consolethree}

[Hint: See previous example.] Now modify your code so that it
produces this output:

\begin{consolethree}
* *

* *

*

* *

* * 
\end{consolethree}

Of course that tells you that you can use a 2D array to do ASCII art!!! You can think of the 2D array as a drawing canvas. There is a big difference between the ASCII art problems in the notes for loops and using a 2D array. You can draw a star any where you like in the array. On the other hand, ASCII art using only loops and print statements requires you to print on the console window from left-to-right, row-by-row.

There is however a restriction when using arrays for ASCII art: You have to specify the row and column sizes of your array. This is what you can do: You can specify a really huge 2D array (say 1000-by-1000) and then only use part of it for drawing.

\sectionthree{Functions}

Everything is similar to functions for 1-dimensional arrays except for one thing. \textbf{LISTEN UP!!!} When you wrote function prototypes or function headers to be general, instead of writing

\verb!void f(int x[3]);!

we write

\verb!void f(int x[]);!

When declaring a 2D array (of ints or doubles or bools or ...) you \textbf{must specify the size of the second dimension}. For instance this is \textbf{WRONG}

\verb!void f(int x[][]);!

This is \textbf{CORRECT}:

\verb!void f(int x[][42]);!

Here's a simple example:

\begin{consolethree}
#include <iostream>

const int COL_SIZE = 2;

void print(int x[][COL_SIZE], int numRows)
{
     for (int row = 0; row < numRows; row++)
     {
          for (int col = 0; col < COL_SIZE; col++)
          {
            std::cout << x[row][col] << ' ';
          }
     }
     std::cout << std::endl; // goto next line
}

int main()
{
    int x[2][COL_SIZE];

    x[0][0] = 1;
    
    x[0][1] = 2;

    x[1][0] = 3;

    x[1][1] = 4;

    print(x, 2);

    return 0;
}
\end{consolethree}

Note that the code now seems to be highly ``nonsymmetric'': you do not have control over the second dimension.

To make the code more symmetric, you might want to do this, especially if you do not need full generality for the first dimension:

\begin{consolethree}
#include <iostream>

const int ROW_SIZE = 2;

const int COL_SIZE = 2;

void print(int x[ROW_SIZE][COL_SIZE])
{
     for (int row = 0; row < ROW_SIZE; row++)
     {
          for (int col = 0; col < COL_SIZE; col++)
          {
               std::cout << x[row][col] << ' ';
          }
     }
     std::cout << std::endl; // goto next line
}

int main()
{
    int x[ROW_SIZE][COL_SIZE];
    x[0][0] = 1;
    x[0][1] = 2;
    x[1][0] = 3;
    x[1][1] = 4;
    print(x);
    return 0;
}
\end{consolethree}

\begin{ex}
Can you change values in a 2D array through a function? Write a simple program and verify.

``Why do I need to specify the size for the second dimension?''

Why? You have to wait till CISS3 ... oh well ... OK I'll explain.

Suppose you have this code:

\begin{consolethree}
int x[4][10];

std::cin << x[2][6];
\end{consolethree}

First of all, our mental picture of the values of array x looks like this:

\begin{align*}
& & & & & & & & & \\
& & & & & & & & & \\
& & & & & & & & & \\
& & & & & & & & & \\
& & & & & & & & & \\
\end{align*}

(i.e., \texttt{x} has 4 rows of 10 integers).

Actually in your computer (at least up to this point in history!), memory is organized in a \textbf{linear} fashion, i.e., \textbf{on a straight line}. So here's a picture of the memory with the part occupied by x shown. I've colored alternate blocks of 10 integers. The location in the computer's memory for

\texttt{x[2][6]} is shown.

\begin{consolethree}[escapeinside=||]
& & & & & & & & & & & & & & & & & & & & & & & & & & & & & & & & & & & & & & & \\
& & & & & & & & & & & & & & & & & & & & & & & & & & & & & & & & & & & & & & & \\
\end{consolethree}

How does C++ locate the value in the array that you want? C++ actually only remembers the location of \texttt{x[0][0]}. When you want \texttt{x[2][6]}, C++ computes the position of \texttt{x[2][6]} as an offset from \texttt{x[0][0]}.

Now the location of \texttt{x[2][6]} is

\begin{center}
2 * 10 + 6
\end{center}

away from the first value \texttt{x[0][0]}. Right? Once C++ knows the computer memory location of the value you want, it can then retrieve the value or assign a value to that spot.

So ... one more time:

\begin{itemize}
\tightlist
\item
  Values are stored in memory which is organized in a straight line.
\item
  For an array, C++ remembers the position of the first element of the array.
\item
  If your program needs \texttt{x[i][j]} of an array x, C++ will compute the location of that value by computing it's offset from \texttt{x[0][0]}.
\end{itemize}

OK. So now look at the above example:

\begin{consolethree}
int x[20][10];

std::cin << x[2][6];
\end{consolethree}

If you need \texttt{x[2][6]}, C++ will look at the value that is

\begin{center}
2 * 10 + 6
\end{center}

away from \texttt{x[0][0]}. But ... notice something? The offset is

\begin{center}
2 * 10 + 6
\end{center}

which depends on 2 and 6 which is provided in your code when you write \texttt{x[2][6]}. But notice that it also depends on ... \textbf{the number 10:}
\begin{center}
2 * \textbf{10} + 6

Where does that come from? It's the \textbf{size of the second dimension}

This is the reason why the size of the second dimension must be stated. Otherwise C++ does not know how to access a required value in the 2D array.

\sectionthree{Exercise: tic-tac-toe}

Write a tic-tac-toe program using a 2D array for the board.

\begin{ex}
Write a function

\begin{consolethree}
const int n = 3;

void init(char board[n][n]); 
\end{consolethree}
that initialization a tic-tao-toe board.. Note your program should have
a global constant for the board size. That way if we change 3 to another
size such as 5, you'll get a 5-by-5 tic-tac-toe board. So in your code
do not assume the board size is 3-by-3.
\end{ex}
